
\documentclass[10pt]{mypackage}

\usepackage[uprightgreeks,noDcommand]{kpfonts}
\renewcommand{\coloneq}{\coloneqq}
\renewcommand{\eqcolon}{\eqqcolon}
%\renewcommand{\mathbb}{\mathds}
%\usepackage{homework}
\usepackage{notes}

\usepackage[ backend=bibtex, style = alphabetic, sorting=ynt ]{biblatex}
\addbibresource{exam_prep_references.bib}

\usepackage{parskip}

\fancyhf{}
\fancyhead[R]{Avinash Iyer}
\fancyhead[L]{Oral Exam Prep: Geometric Group Theory}
\fancyfoot[C]{\thepage}

\setcounter{secnumdepth}{0}

\begin{document}
\RaggedRight
This is one of my two preparation documents for my oral exam, focused on geometric/combinatorial group theory.
\section{Preliminaries}%
\subsection{Group-Theoretic Constructions}%
There are two major ways we may construct new groups out of old: extensions/semidirect products and pushouts/amalgamated free products. We will provide a brief overview of both of these.
\begin{definition}
  Let $Q$ and $N$ be groups. We say a group $G$ is an extension of $Q$ and $N$ if there is an exact sequence
  \begin{align*}
    1 \rightarrow N\rightarrow G\rightarrow Q\rightarrow 1.
  \end{align*}
  If the extension splits --- i.e., there is a section $Q\rightarrow N$ --- then $G\cong Q\ltimes_{\varphi}N$ is a semidirect product for some $\varphi\colon Q\rightarrow \aut(N)$.
\end{definition}
\begin{example}\hfill
  \begin{itemize}
    \item If the homomorphism $\varphi\colon Q\rightarrow \aut(N)$ is trivial, then $N\rtimes_{\varphi}Q\cong N\times Q$ is a direct product.
    \item The map $\Z/2\Z\rightarrow \aut\left( \Z/n\Z \right)$ given by inversion yields the dihedral group
      \begin{align*}
        D_n &\cong \Z/2\Z\ltimes \Z/n\Z.
      \end{align*}
    \item If $G$ is any group, the lamplighter group of $G$ is the semidirect product $\Z\ltimes G^{\Z}$ with homomorphism $\varphi\colon \Z\rightarrow \aut\left( G^{\Z} \right)$ given by
      \begin{align*}
        \varphi(z) &= \left( \left( g_n \right)_{n\in \Z}\mapsto \left( g_{n + z} \right)_{n\in \Z} \right).
      \end{align*}
    \item the wreath product of $G$ and $H$, denoted $G\wr H$, is the semidirect product $H\ltimes G^{H}$ with $\varphi\colon H\rightarrow \aut\left( G^{H} \right)$ given by the shift action
      \begin{align*}
        \varphi(h) &= \left( \left( g_k \right)_{k\in H}\mapsto \left( g_{kh} \right)_{k\in H} \right).
      \end{align*}
  \end{itemize}
\end{example}
\begin{definition}
  Let $A$ be a group, $\alpha_1\colon A\rightarrow G_1$, and $\alpha_2\colon A\rightarrow G_2$ group homomorphisms. We say a group $G$ is a \textit{pushout} of $G_1$ and $G_2$ over $A$ with respect to $\alpha_1$ and $\alpha_2$ if there are homomorphisms $\beta_1\colon G_1\rightarrow G$ and $\beta_2\colon G_2\rightarrow G$ satisfying $\beta_1\circ \alpha_1 = \beta_2\circ \alpha_2$ and the universal property for pushouts.

  That is, for every group $H$ and all group homomorphisms $\varphi_1\colon G_1\rightarrow H$ and $\varphi_2\colon G_2\rightarrow H$ satisfying $\varphi_1\circ \alpha_1 = \varphi_2\circ \alpha_2$, there is exactly one homomorphism $\varphi\colon G\rightarrow H$ satisfying $\varphi\circ\beta_i = \varphi_i$ for each $i=1,2$. We denote the pushout by $G_1\ast_{A}G_2$.
  \begin{itemize}
    \item If $A = \set{e}$ is the trivial group, then the pushout $G_1\ast_A G_2 = G_1\ast G_2$ is called the \textit{free product} of $G_1$ and $G_2$.
    \item If $\alpha_1$ and $\alpha_2$ are injective, then $G_1\ast_{A}G_2$ is called an amalgamated free product of $G_1$ and $G_2$ over $A$.
  \end{itemize}
\end{definition}
\begin{example}\hfill
  \begin{itemize}
    \item The free group $F_n$ on $n$ letters is the $n$-fold free product $\Z\ast\cdots\ast\Z$.
    \item The infinite dihedral group $D_{\infty}\cong \Z/2\Z\ltimes \Z$ is isomorphic to the free product $\Z/2\Z\ast \Z/2\Z$.
    \item The matrix group $\SL_2\left(\Z\right)$ is isomorphic to the amalgamated free product $\Z/6\Z\ast_{\Z/2\Z}\Z/4\Z$.
    \item If $X$ and $Y$ are pointed, non-empty, path-connected topological spaces with nonempty and path-connected intersection, then the fundamental group of their union (which is a type of pushout) is given by the amalgamated free product of their fundamental groups over the fundamental group of the intersection.
  \end{itemize}
\end{example}
\begin{theorem}
  Pushout groups exist and are unique up to isomorphism.
\end{theorem}
\begin{proof}
  Take presentations
  \begin{align*}
    G_i &= \left\langle \set{x_g : g\in G_i}|R_i \right\rangle.
  \end{align*}
  Then, the presentation
  \begin{align*}
    G\coloneq \left\langle \set{x_g : g\in G_1}\cup \set{x_g : g\in G_2} | \set{x_{\alpha_1(a)}x_{\alpha_2(a)}^{-1} : a\in A} \cup R_1 \cup R_2 \right\rangle.
  \end{align*}
  The group homomorphisms $\beta_i\colon G_i\rightarrow G$ taking $g\mapsto x_g$, along with the relations $\set{x_{\alpha_1(a)}x_{\alpha_2(a)}^{-1} : a\in A}$, give $\beta_1\circ \alpha_1 = \beta_2\circ \alpha_2$. The universal property then follows from the universal property of generators and relations.
\end{proof}
\begin{definition}
  Let $G$ be a group, and let $A,B\subseteq G$ be isomorphic subgroups with isomorphism $\theta\colon A\rightarrow B$. The \textit{HNN extension} of $G$ with respect to $\theta$ is the group
  \begin{align*}
    G \ast_{\theta} &\coloneq \left\langle \set{x_g : g\in G}\cup\set{t} | \set{t^{-1}x_a t x_{\theta(a)}^{-1} : a\in A}\cup R \right\rangle,
  \end{align*}
  where $R$ is the relations for $G$. We say that $t$ is the stable letter of the HNN extension.
\end{definition}
The HNN extension is the universal (in a sense) way to force isomorphic subgroups to be conjugate to each other.
\subsection{Group Homology/Cohomology}%
\begin{definition}
  Let $R$ be a commutative unital ring. The \textit{group ring} is the ring $R[G]$ with
  \begin{itemize}
    \item underlying abelian group $\bigoplus_{g\in G}Rg$, where the $g$ are formal symbols corresponding to the elements of $G$;
    \item multiplication given by the $R$-bilinear extension of the group multiplication $\left( g,h \right)\mapsto gh$.
  \end{itemize}
\end{definition}
If $G$ is a non-abelian group and $R$ is non-trivial, then the group ring $R[G]$ is not commutative.
\begin{definition}
  Let $G$ be a group.
  \begin{itemize}
    \item The \textit{simplicial resolution} of $G$, denoted $C_{\ast}(G)$, is the chain complex of $\Z$-modules whose chains are defined by
      \begin{align*}
        C_n &\coloneq \bigoplus_{G^{n+1}} \Z,
      \end{align*}
      with boundary maps given by $\partial_n\colon C_n(G)\rightarrow C_{n-1}(G)$ taking
      \begin{align*}
        \partial_n\left( \sum_{g\in G^{n+1}}a_g \left( g_0,\dots,g_n \right) \right) &= \sum_{g\in G^n} a_g\sum_{j=0}^{n} \left( -1 \right)^{j} \left( g_0,\dots,g_{j-1},g_{j+1},\dots,g_n \right).
      \end{align*}
    \item If $A$ is a left $\Z[G]$-module, then the simplicial chain complex of $G$ with $A$-coefficients is given by $C_n\left( G;A \right) = A\otimes_{\Z[G]}C_{\ast}(G)$ with boundary operator $\partial_{\ast}^A = \id_A\otimes_{\Z[G]}\partial_{\ast}$. The elements of $\ker\left( \partial_n^A \right)$ are called the $n$-cycles.

      The $n$-th homology of $G$ with $A$-coefficients is the $\Z$-module given by
      \begin{align*}
        H_n\left( G;A \right) &\coloneq \frac{\ker\left( \partial_n^A\colon C_n\left( G;A \right)\rightarrow C_{n-1}\left( G;A \right) \right)}{\img\left( \partial_{n+1}^{A}\colon C_{n+1}\left( G;A \right)\rightarrow C_n\left( G;A \right) \right)}.
      \end{align*}
    \item Let $A$ be a left $\Z[G]$-module. The simplicial cochain complex of $G$ with $A$-coefficients is defined by
      \begin{align*}
        C^{\ast}\left( G;A \right) \coloneq \Hom_{\Z[G]} \left( C_{\ast}(G),A \right)
      \end{align*}
      with coboundary operator $\partial_A^{\ast}$ given by the $A$-dual of $\partial_{\ast}$. Elements of $\ker\left( \delta_A^n \right)$ are called the $n$-cocycles. The $n$th cohomology of $G$ with $A$-coefficients is the $\Z$-module
      \begin{align*}
        H^n\left( G;A \right) &\coloneq \frac{\ker\left( \partial^n_A\colon C^n\left( G;A \right)\rightarrow C^{n+1}\left( G;A \right) \right)}{\img\left( \partial^{n-1}_A\colon C^{n-1}\left( G;A \right)\rightarrow C^n\left( G;A \right) \right)}.
      \end{align*}
  \end{itemize}
\end{definition}
\subsection{Nilpotent Groups}%
\begin{lemma}
  Let $G$ be a group, let $N\trianglelefteq G$ be a normal subgroup, and let $H,K,L\leq G$ be subgroups. If $\left[ \left[ H,K \right],L \right]$ and $ \left[ \left[ K,L \right],H \right] $ are contained in $N$, then so is $ \left[ \left[ L,H \right],K \right] $.
\end{lemma}
\begin{proof}
  The subgroups $ \left[ \left[ H,K \right],L \right] $, $ \left[ \left[ K,L \right],H \right] $, and $ \left[ \left[ L,H \right],K \right] $ are generated by the commutators $ \left[ h,k,\ell \right] $, $ \left[ k,\ell,h \right] $, and $ \left[ \ell,h,k \right] $ respectively, where $h\in H$, $k\in K$, and $\ell\in L$. The Hall--Witt identity says that
  \begin{align*}
    e &= b^{-1}\left[ a,b^{-1},c \right]b\\
      &=  c^{-1} \left[ b,c^{-1},a \right]c\\
      &= a^{-1} \left[ c,a^{-1}b \right]a,
  \end{align*}
  so all of these commutators are contained in $N$.
\end{proof}
\begin{definition}
  Let $G$ be a group. The \textit{lower central series} of $G$ is the sequence $ \left( \gamma_k(G) \right)_{k\geq 1} $ of subgroups defined recursively by
  \begin{align*}
    \gamma_1(G) &\coloneq G\\
    \gamma_{k+1}(G) &\coloneq \left[ \gamma_k(G),G \right].
  \end{align*}
  We say a group $G$ is \textit{nilpotent} if its lower central series is finite. That is, if there is some integer $k\geq 1$ such that $\gamma_k(G) = e$. The smallest such $k$ is called the \textit{nilpotence class} (or nilpotency class) of $G$.
\end{definition}
Inductively, we observe that $\gamma_k(G)$ is characteristic in $G$ with $\gamma_{k+1}(G)\subseteq \gamma_k(G)$, meaning that the lower central series is a descending normal series.
\begin{lemma}
  Let $G$ be a group. Then,
  \begin{align*}
    \left[ \gamma_i(G),\gamma_j(G) \right] &\subseteq \gamma_{i+j}(G).
  \end{align*}
\end{lemma}
\begin{proof}
  We prove by induction on $i$. By the definition of $\gamma_{i+1}$, we have that $ \left[ \gamma_1(G),\gamma_j(G) \right] = \left[ G,\gamma_j(G) \right] = \gamma_{j+1}(G) $.

  Now, we show that
  \begin{align*}
    \left[ \gamma_{i+1}(G),\gamma_j(G) \right] &\subseteq \gamma_{i + j + 1}(G).
  \end{align*}
  First, we claim that given $a,b\in \gamma_{i+1}(G)$ and $c\in \gamma_j(G)$, we have
  \begin{align*}
    \left[ ab,c \right] &= \left[ a,c \right]\left[ b,c \right]\text{ mod }\gamma_{i + j + 1}(G)\\
    \left[ a^{-1},c \right] &= \left[ a,c \right]^{-1}\text{ mod }\gamma_{i + j + 1}(G).
  \end{align*}
  Since $a\in \gamma_{i+1}(G)\subseteq \gamma_i(G)$, the induction hypothesis has $\left[ a,c \right]\in \left[ \gamma_{i}(G),\gamma_j(G) \right]\subseteq \left[ \gamma_{i+j}(G) \right]$. Consequently, $\left[ \left[ a,c \right],b \right]\in \left[ \gamma_{i+j}(G),\gamma_j(G) \right]\subseteq \left[ \gamma_{i+j}(G),G \right] = \gamma_{i + j + 1}(G)$.

  By standard commutator identities, we have
  \begin{align*}
    \left[ ab,c \right] &= \left[ a,c \right]\left[ a,c,b \right]\left[ b,c \right]\\
                        &= \left[ a,c \right]\left[ b,c \right]\left( \left[ b,c \right]^{-1}\left[ a,c,b \right] \left[ b,c \right]\right).
  \end{align*}
  Since $\gamma_{i + j + 1}(G)$ is normal, it follows that $\left[ b,c \right]^{-1}\left[ a,c,b \right]\left[ b,c \right]\in \gamma_{i + j + 1}(G)$, so this gives our first desired result. Similarly, by taking $b = a^{-1}$, we obtain the second desired result.

  Therefore, we only need to show that $\left[ a,x,b \right]\in \gamma_{i + j + 1}(G)$ For all $a\in \gamma_i(G)$, $x\in G$, and $b\in \gamma_j(G)$.

  From the Hall--Witt identity, we may write
  \begin{align*}
    \left[ a,x,b \right] &= x^{-1} \left( \left( a^{-1} \left[ b,a^{-1}x^{-1} \right]a \right)^{-1}\left( b^{-1}\left[ x^{-1},b^{-1}a \right]b \right)^{-1} \right)x.
  \end{align*}
  while by the inductive hypothesis, $ \left[ b,a^{-1},x^{-1} \right]\in \left[ \gamma_{i+j}(G),G \right] = \gamma_{i+j+1}(G) $, and $\left[ x^{-1},b^{-1},a \right]\in \left[ \gamma_{j+1}(G),\gamma_i(G) \right]\subseteq \gamma_{i + j + 1}(G)$, seeing as $\gamma_{i + j + 1}(G)$ is a normal subgroup (hence the entire conjugacy class must land in it if any conjugate lands in it).
\end{proof}
We say a descending series $G = G_1\geq G_2\geq\cdots\geq G_k\geq G_{k+1}\geq\cdots$ is \textit{central} if $ \left[ G_i,G_j \right]\subseteq G_{i + j} $ for all $i,j\geq 1$. This condition implies that each $G_i$ is normal in $G$.
\begin{lemma}
  Let $G$ be a group. Suppose $\left( G_k \right)_{k\geq 1}$ is a descending central series for $G$. Then, $G_i\supseteq \gamma_i(G)$ for each $i\geq 1$.
\end{lemma}
\begin{proof}
  We proceed by induction. First, $G_1 = G = \gamma_1(G)$.

  Suppose the statement holds for $i$. The defining property of a central series and the inductive hypothesis give
  \begin{align*}
    G_{i+1} &\supseteq \left[ G_i,G \right]\\
    &\supseteq \left[ \gamma_i(G),G \right]\\
    &= \gamma_{i+1}(G).
  \end{align*}
\end{proof}
In particular, this means that a group $G$ is nilpotent if and only if it admits any finite descending central series.
\begin{proposition}
  Let $G$ be a nilpotent group with nilpotence class $c$. Then, any subgroup and any quotient of $G$ is nilpotent of nilpotence class at most $c$.
\end{proposition}
\begin{proof}
  If $H\leq G$ is a subgroup, inductively we have $\gamma_i(H)\subseteq \gamma_i(G)$.

  If $\varphi\colon G\rightarrow \overline{G}$ is a surjective homomorphism, then we have $\varphi\left( \left[ H,K \right] \right) = \left[ \varphi\left( H \right),\varphi\left( K \right) \right]$ for all subgroups $H,K\leq G$. In particular, this means that inductively we have $\gamma_{i+1}\left( \varphi(G) \right) = \varphi\left( \gamma_{i+1}(G) \right)$.
\end{proof}
\begin{definition}
  Let $G$ be a group. The \textit{upper central series} of $G$ is the ascending series $\left( Z_k(G) \right)_{k\geq 0}$, where
  \begin{align*}
    Z_0(G) &= e\\
    Z_{i+1}(G) &= \pi_i^{-1}\left( Z\left( G/Z_i(G) \right) \right),
  \end{align*}
  where $\pi_i\colon G\rightarrow G/Z_i(G)$ is the projection homomorphism, and $Z(H)$ is the center of a group $H$.
\end{definition}
In particular, we have $Z\left( G/Z_i(G) \right) = Z_{i+1}(G)/Z_i(G)$.

If $G$ admits a finite upper central series, then the descending central series $\left( Z_{k-i+1} \right)_{i=1}^{k+1}$ is central.
\begin{lemma}
  Let $G$ be a group. Suppose $G$ admits a finite central series $G = G_1\geq G_2\geq\cdots\geq G_r = \set{e_G}$. Then, $Z_i(G)\supseteq G_{r-i}$ for all $i\geq 0$. In particular, this means $Z_{r-1}(G) = G$.
\end{lemma}
\begin{proof}
  We induct on $i$. First, observe that $Z_0(G) = \set{e_G} = G_r$.

  We will show that $Z_{i+1}(G)\supseteq G_{r-i-1}$. That is, we have to show that $\pi_i\left( G_{r-i-1} \right)\subseteq Z\left( G/Z_i(G) \right)$. This is equivalent to showing that $\left[ G_{r-i-1},G \right]\subseteq Z_i(G)$. Yet, this follows from the fact that the sequence $\left( G_i \right)_{i\geq 0}$ is central, whence $\left[ G_{i-r-1},G \right]\subseteq G_{r-i}$, and $G_{r-i}\subseteq Z_i(G)$ by the induction hypothesis.
\end{proof}
In particular, observe that $\gamma_i(G)\subseteq Z_{c-i}(G)$ if $G$ is nilpotent of nilpotence class $c$.
\subsubsection{Finitely Generated Nilpotent Groups}%
We start by considering torsion-free finitely generated nilpotent groups. That is, finitely generated nilpotent groups where no nontrivial element has finite order.
\begin{lemma}
  Let $G$ be a torsion free (not necessarily finitely generated) nilpotent group with upper central series $\left( Z_k(G) \right)_{k\geq 0}$. Then, all the factors $Z_{i+1}(G)/Z_i(G)$ are torsion-free abelian groups.
\end{lemma}
\begin{proof}
  We use induction on nilpotence class. If $c = 1$, then $Z_1(G) / Z_0(G) = G$ is abelian, so it is a torsion-free abelian group by assumption.

  Suppose it holds true for all nilpotent groups of nilpotence class $\left( c-1 \right)$. Let $G$ have nilpotence class $c$. Set $H = G/Z_1(G)$. Then, we have $Z_i(H) = Z_{i+1}(G)/Z_i(G)$ by induction and the definition of the upper central series. In particular, $H$ has nilpotence class $\left( c-1 \right)$.

  By the third isomorphism theorem, we have $Z_{i+1}(G)/Z_i(G)\cong \left( Z_{i+1}(G)/Z_1(G) \right)/\left( Z_i(G)/Z_1(G) \right) \cong Z_i(H)/Z_{i-1}(H)$, so it suffices to show that $H$ is torsion-free to complete the proof.

  Let $a\in G$. Suppose there exists an integer $m\geq 1$ such that $a^{m}\in Z_1(G)$, and suppose toward contradiction that $a\notin Z_1(G)$. Notice that $Z_1(G) = Z(G)$, so by assumption there is some $b\in G$ such that $\left[ a,b \right]\neq e$. Let $1\leq i \leq c-1$ be the maximal index for which there exists some $b\in \gamma_i(G)$ such that $\left[ a,b \right]\neq e$, where $\left( \gamma_i(G) \right)_{i=1}^{c+1}$ is the lower central series. By the various commutator identities, we have
  \begin{align*}
    e &= \left[ a^{m},b \right]\\
      &= \left[ aa^{m-1},b \right]\\
      &= a^{1-m} \left[ a,b \right]a^{m-1} \left[ a^{m-1},b \right].
  \end{align*}
  Since $\left[ a,b \right]\in \gamma_{i+1}(G)$, we must have that $a$ commutes with $\left[ a,b \right]$, so $a^{1-m}\left[ a,b \right]a^{m-1} = \left[ a,b \right]$, and we have $e = \left[ a^{m},b \right] = \left[ a,b \right]\left[ a^{m-1}b \right]$. Iterating this process gives $e = \left[ a,b \right]^{m}$, but since $G$ is torsion free, this gives that $\left[ a,b \right] = e$, which is a contradiction of the definition of $b$. Thus, $a\in Z_1(G)$, hence $H = G/Z_1(G)$ is torsion-free.
\end{proof}

\section{Groups and Graphs}%
One of the first connections between groups and geometry is the notion of a Cayley graph, and the way that the action of a group on a graph is connected to the structure of the group.
\subsection{Graph Theory Review}%
\begin{definition}\hfill
  \begin{itemize}
    \item A graph is a pair $X = \left( V,E \right)$, where $E\subseteq V^{[2]} \coloneq \set{e : e\subseteq V,\left\vert e \right\vert = 2}$. Elements of $E$ are denoted by $\set{v_1,v_2}$ if the graph is undirected. If the graph is directed, then elements of $E$ are denoted by tuples $\left( v_1,v_2 \right)$.
    \item If $\left( V,E \right)$ is a graph, then we say vertices $v,v'\in V$ are adjacent if $\set{v,v'}\in E$. The number of neighbors is the degree of the vertex.
    \item If $X = \left( V,E \right)$ and $X' = \left( V',E' \right)$ are graphs, then we say $X$ and $X'$ are isomorphic if there is a graph isomorphism between $X$ and $X'$ --- that is, a bijection $f\colon V\rightarrow V'$ such that for any $v,w\in V$, we have $\set{v,w}\in E$ if and only if $\set{f(v),f(w)}\in E'$.
  \end{itemize}
\end{definition}
\begin{definition}
  Let $X = \left( V,E \right)$ be a graph. 
  \begin{itemize}
    \item A path in $X$ is a sequence $\left[ v_0,\dots,v_n \right]$ of distinct vertices such that $\set{v_j,v_{j+1}}\in E$.
    \item A graph is called connected if every pair of its vertices can be connected by a path in $X$.
    \item A cycle of length $n$ is a path $\left[ v_0,\dots,v_{n} \right]$ such that $\set{v_n,v_0}\in E$.
    \item A tree is a connected graph without cycles. Note that this is equivalent to a graph where every pair of vertices has exactly one path connecting them.
    \item If $X$ is a connected graph, a \textit{spanning tree} for $X$ is a tree that contains all the vertices of $X$.
  \end{itemize}
\end{definition}
In algebraic topology, we use spanning trees as a way to compute fundamental groups. This is done by using the fact that spanning trees are CW subcomplexes (viewing the graph as a CW complex).
\subsection{Cayley Graphs}%
\begin{definition}
  Let $G$ be a group, and let $S\subseteq G$ be a generating set for $G$ without the identity. The Cayley graph of $G$ with respect to $S$ is the graph $ \operatorname{Cay}\left( G,S \right) $, with
  \begin{itemize}
    \item vertices given by elements of $G$;
    \item edges given by
      \begin{align*}
        \set{\set{g,gs} : g\in G,s\in S\cup S^{-1}}.
      \end{align*}
  \end{itemize}
\end{definition}
\begin{example}\hfill
  \begin{itemize}
    \item The Cayley graph of $\Z$ with respect to $\set{1}$ appears as an infinite line, while the Cayley graph of $\Z$ with respect to the generating set $\set{2,3}$ appears as a particular $4$-regular graph. See \cite[Figure 3.3]{loh_geometric_group_theory}.
    \item The Cayley graph of $\Z^2\subseteq \R^2$ with respect to the generating set $\set{\left( 1,0 \right),\left( 0,1 \right)}$ appears as the integer lattice.
    \item The Cayley graph of $\Z/6\Z$ with respect to $\set{[1]}$ is a cycle graph.
    \item The Cayley graph of the symmetric group $S_3$ with respect to the generating set $\set{\sigma,\tau}$, where we have
      \begin{align*}
        S_3 &\coloneq \left\langle \sigma,\tau | \sigma^3,\tau^2,\tau\sigma\tau\sigma \right\rangle.
      \end{align*}
    \item The Cayley graph of any group with respect to the generating set given by the group itself is a complete graph.
  \end{itemize}
\end{example}
One of the seemingly simple questions one might ask is when exactly finitely generated groups $G$ and $H$ with finite generating sets $S\subseteq G$ and $T\subseteq H$ are such that $ \operatorname{Cay}\left( G,S \right) \cong \operatorname{Cay}\left( H,T \right) $ are isomorphic as graphs. This is essentially a question about Cayley graphs whose automorphisms are affine (i.e., given by a group isomorphism and translation by a fixed group element).
\begin{itemize}
  \item Cayley graphs of finitely generated abelian groups satisfy the rigidity question; the automorphisms are affine on the free part, and Cayley graphs remember the rank of the group and the size of the torsion subgroup.
  \item Automorphisms of Cayley graphs of finitely generated torsion-free nilpotent groups are affine.
  \item Finitely generated free groups admit isomorphic Cayley graphs if and only if they have the same rank.
\end{itemize}
Given a presentation of a group, the presentation complex is a two-dimensional CW complex given by taking one zero-cell, adding circles for every generator, and attaching disks for every relation (attached by the word for the relation).

The presentation complex associated to $ \left\langle x,y | xyx^{-1}y^{-1} \right\rangle $ is the $2$-torus. The presentation complex associated to $ \left\langle x|x^2 \right\rangle $ is $ \mathbb{RP}^2 $.

In general, any group admits a classifying space (known as the Eilenberg--MacLane space) whose fundamental group is the underlying group and whose higher homotopy groups are trivial. These spaces are unique up to homotopy equivalence. For example, the torus is a classifying space for $\Z^2$ and $ \mathbb{RP}^{\infty} $ is a classifying space for $\Z/2\Z$. The $1$-skeleton of the universal cover of the presentation complex of the presentation $ \left\langle X|R \right\rangle $ is basically the Cayley graph of the underlying group (save for generators of order $2$).
\begin{theorem}
  Let $F$ be a free group freely generated by $S\subseteq F$. Then, $ \operatorname{Cay}\left( F,S \right) $ is a tree.
\end{theorem}
\begin{proof}
  Every element of $F$ is a unique word in elements of $S$, so there is a unique path between the identity and any word, so there is a unique path between any two elements, so $ \operatorname{Cay}\left( F,S \right) $ is a tree.
\end{proof}
There is a converse, but it requires a caveat. In particular, we see that $ \operatorname{Cay}\left( \Z/2\Z,[1] \right) $ is a one-edge two-vertex graph (which is a tree). 
\begin{theorem}
  Let $G$ be a group, $S\subseteq G$ a generating set satisfying $st\neq e$ for all $s,t\in S$. If $\operatorname{Cay}\left( G,S \right)$ is a tree, then $G \cong F(S)$ is freely generated by $S$.
\end{theorem}
\begin{proof}
  Let $S$ and $G$ be such that $\operatorname{Cay}\left( G,S \right)$ is a tree. We will show that $G$ is isomorphic to $F(S)$ by an isomorphism that is the identity on $S$.

  First, we observe that, by the universal property applied to the inclusion $S\subseteq G$, there is a group homomorphism $\varphi\colon F(S)\rightarrow G$ that is the identity on $S$. Since $S$ generates $G$, it follows that $\varphi$ is surjective.

  Assume toward contradiction that $\varphi$ is not injective. Let $s_1\cdots s_n\in F(S)\setminus \set{\ve}$ be a word of minimal length mapped to $e\in G$ by $\varphi$.

  First, since $\varphi|_{S} = \id_S$, we must have $n > 1$.

  Next, if $n = 2$, then we would have
  \begin{align*}
    e &= \varphi\left( s_1s_2 \right)\\
      &= \varphi\left( s_1 \right)\varphi\left( s_2 \right)\\
      &= s_1s_2,
  \end{align*}
  contradicting the fact that $s_1\cdots s_n$ is reduced with $st\neq e$ in $G$.

  Next, suppose $n\geq 3$. Inductively, we may consider the sequence of elements $g_0,\dots,g_{n-1}$ in $G$ given by $g_{j+1} = g_js_{j+1}$ for all $0\leq j \leq n-2$. By minimality of the word $s_1\cdots s_n$, the elements $g_0,\dots,g_{n-1}$ are distinct, and since $g_n = s_1\cdots s_n = e$, it follows that $\left[ g_0,\dots,g_{n-1} \right]$ is a cycle in $\operatorname{Cay}\left( G,S \right)$. Yet, this contradicts the fact that $\operatorname{Cay}\left( G,S \right)$ is a tree. Thus, $\varphi$ is injective.
\end{proof}
It turns out that, generally speaking, we are able to reconstruct a Cayley graph from any action of a group on a graph.
\begin{proposition}
  Let $G$ be a group, and let $G$ act on a connected graph $X = \left( V,E \right)$ by graph automorphisms. If this action is free and transitive on the set $V$ of vertices of $X$, then for a distinguished vertex $x$, the set
  \begin{align*}
    S \coloneq \set{s\in G : \set{x,s.x}\in E}
  \end{align*}
  generates $G$, and $\operatorname{Cay}\left( G,S \right)\cong X$ are isomorphic as graphs.
\end{proposition}
\begin{proof}
  Since the action of $G$ on the vertices is free and transitive, we may identify $V$ with $G$ by the identification $g\mapsto g.x$.

  We start by showing that $S$ generates $G$. Let $g\in G$. Since the graph $X$ is connected, there is a path $\left[ x_0,x_1,\dots,x_n \right]$, where $x_0 = x$ and $x_n = g.x$, that joins $x$ and $g.x$. We will translate this path into steps in the group $G$ as follows.

  For each $0\leq j \leq \left( n-1 \right)$, define
  \begin{align*}
    g_j &= \varphi^{-1}\left( x_j \right)
  \end{align*}
  in $G$. Then, we take $s_j = g_j^{-1}g_{j+1}$. We will show that $s_j\in S$. Since $\left[ x_0,\dots,x_n \right]$ is a path in $X$, we have $\set{x_j,x_{j+1}}$ is an edge of $X$. Since $G$ acts by graph automorphisms, we have that $\set{g_j^{-1}.x_j,g_j^{-1}.x_{j+1}}$ is also an edge of $X$, but this is equal to $\set{x,\left( g_j^{-1}g_{j+1} \right).x}$. In particular, this means that
  \begin{align*}
    g &= g_n\\
      &= g_0g_0^{-1}g_1g_1^{-1}\cdots g_{n-1}g_{n-1}^{-1}g_n\\
      &= s_0s_1\cdots s_{n-1},
  \end{align*}
  so $S$ is a generating set for $G$.

  We only need to show now that $\operatorname{Cay}\left( G,S \right)$ is isomorphic to the graph $X$. We will show that $\varphi$ induces a graph isomorphism. If $g,h\in G$, then we observe that $\set{\varphi(g),\varphi(h)} = \set{g.x, h.x}$. Since $G$ acts by graph automorphisms on $X$, the set is an edge of $X$ if and only if $\set{x,\left( g^{-1}h \right).x}$ is an edge of $X$, which holds if and only if $g^{-1}h\in S$, which is equivalent to $\set{g,h}$ being an edge of $\operatorname{Cay}\left( G,S \right)$. Thus, $ \operatorname{Cay}\left( G,S \right) $ is isomorphic to $X$.
\end{proof}

\subsection{Actions of Groups on Trees}%
Similar to how a group is a free group with generating set $S$ if and only if $\operatorname{Cay}\left( G,S \right)$ is a tree, there is a converse direction. That is, a group is free if and only if it admits a free action on a non-empty tree.

One direction emerges immediately from the fact that left-translation is a free action of $G$ on itself that acts on its Cayley graph.

The reverse direction is a little more difficult. If $G\curvearrowright T$ is an action, we will have to find a tree that is a Cayley graph for $G$.
\begin{definition}
  Let $G$ be a group that acts on a connected graph $X$ by graph automorphisms. A \textit{spanning tree} of this action is a subgraph of $X$ that is both a tree and contains exactly one vertex of every orbit of the action of $G$ on the vertices of $X$.
\end{definition}
\begin{theorem}
  Every action of a group on a connected graph by automorphisms has a spanning tree.
\end{theorem}
\begin{proof}
  Let $G$ be a group acting on a connected graph $X$. We may assume that $X$ is nonempty.

  Consider the set $ T_G $ of all subtrees of $X$ that contain at most one vertex of every orbit. This is a nonempty set since $X$ has at least one vertex. The set $T_G$ is partially ordered by the subgraph relation, and every totally ordered chain of $T_G$ has an upper bound in $T_G$ (namely, the union of the trees in this chain). Thus, by Zorn's Lemma, there is a maximal element $T$ in $T_G$ that is necessarily nonempty since $X$ is nonempty.

  Suppose toward contradiction that $T$ is not a spanning tree for the action of $G$. Then, there is a vertex $v$ such that none of the vertices of the orbit $G.v$ is a vertex of $T$. We will show that there is some such $v$ that is adjacent to a vertex of $T$.

  Since $X$ is connected, there is a path $p$ connecting some vertex $u$ of $T$ with $v$, which we will label as
  \begin{align*}
    p &= [u,\dots,v',\dots,v],
  \end{align*}
  where $v'$ is the first vertex of $p$ that is not in $T$.
  \begin{enumerate}[(i)]
    \item If none of the vertices of the orbit $G.v'$ are contained in $T$, then $v'$ is the vertex with our desired property.
    \item If there is some $g\in G$ such that $g.v'$ is a vertex of $T$, write a subpath $p' = \left[ v',\dots,v \right]$ of $p$, and consider the path
      \begin{align*}
        g.p' &= \left[ g.v',\dots,g.v \right].
      \end{align*}
      Notice that $g.v$ is not contained in $T$ since none of the elements of the orbit of $v$ are contained in $T$ by assumption. Since the length of $p'$ is smaller than the length of $p$, we may iterate this process to eventually obtain our desired vertex.
  \end{enumerate}
  Having obtained our vertex $v$ with none of the vertices of the orbit $G.v$ in $T$ and some neighbor $u$ of $v$ is contained in $T$, we may add the edge $\set{u,v}$ to $T$ to produce a tree in $T_G$ that contains $T$ as a proper subgraph. This yields the contradiction, as we had assumed $T$ was maximal. Thus, $T$ is a spanning tree for the action of $G$ on $X$.
\end{proof}
In order to show the reverse direction, we will have to determine, given an action $G\curvearrowright T$, a set $S$ such that $\operatorname{Cay}\left( G,S \right)$ is a tree.
\begin{theorem}
  If a group $G$ acts freely on a tree $T$, then $G$ is a free group.
\end{theorem}
\begin{proof}
  Let $T'$ be a spanning tree for the action of $T$. Say two translations $g.T'$ and $h.T'$ are adjacent if $g.T'\neq h.T'$ and there exist vertices $v\in g.T'$ and $w\in h.T'$ such that $\set{v,w}$ is an edge of $T$. Define the graph $X$ by saying $V(X)$ consists of the set $\set{g.T' : g\in G}$ and $E(X)$ consists of all the adjacent trees $\set{g.T',h.T'}$.

  We will start by showing that if $g\neq h$, then $g.T'\neq h.T'$. We may reduce this question to showing that for any $g\neq e$, we have that $g.T'\neq T'$. Yet, this follows directly from our assumptions; there is exactly one representative from each orbit in the spanning tree, and we cannot have $g.v = v$ for any $v\in T'$ unless $g = e$ by the assumption of the freeness of the action $G$.

  For any adjacent $g.T'$ and $h.T'$, we claim that there is a unique connecting edge. We observe that $g.T'$ and $h.T'$ are adjacent if and only if $h^{-1}g.T'$ is adjacent to $T'$ --- that is, there exists $u\in T'$ and $v\in h^{-1}g.T'$ such that $\set{u,v}$ is an edge. Yet, since $u$ is a unique representative of the orbit of $v$, it follows that there is only one such edge. (In Löh's book, such an edge is dubbed an ``essential edge'' on page 89.)

  We observe that, since each copy of $T'$ is a tree (hence connected), and $T$ itself is a tree, it follows that there can only be one unique path between any two vertices in $T'$ as else that implies the existence of two non-unique paths in the original graph $T$.

  Finally, since $V(X) \curvearrowleft G$ freely by taking $\left( h.T' \right).g = \left( hg \right).T'$, we must have that $\operatorname{Cay}\left( G,S \right)\cong X$. Since $X$ is a tree, it follows that $G$ is a free group.
\end{proof}
The following is a more topological proof that relies heavily on covering space theory.
\begin{proof}[Proof via Covering Spaces]
  Let $G\curvearrowright T$ freely by graph automorphisms. We may view $T$ as a one-dimensional CW complex; since $T$ is a tree, this one-dimensional CW complex is simply connected, so the covering map $T\rightarrow T/G$ has it such that $T$ is the universal cover for the CW complex $T/G$. In particular, we have that the deck transformations of $T$ are exactly $G$, and $\pi_1\left( T/G \right)\cong G$.

  Since any one-dimensional CW complex is homotopy-equivalent to a bouquet of circles, the Seifert--Van Kampen theorem gives that $\pi_1\left( T/G \right)\cong F(S)$ for some generating set $S$, so that $G$ is free.
\end{proof}
We now focus on some applications. 
\begin{theorem}[Nielsen--Schreier Theorem]
  Let $G$ be a free group, and let $H\leq G$ be a subgroup. Then, $H$ is a free group.
\end{theorem}
\begin{proof}
  If $G\curvearrowright T$ acts freely, then so too does $H$. 
\end{proof}
\subsection{Ping-Pong Lemma}%
In order to determine the existence of free subgroups of some group $G$, we use the much celebrated ping-pong lemma.
\begin{theorem}[Ping-Pong Lemma]
  Let $G_1$ and $G_2$ be subgroups of a group $G$ with $\left\vert G_1 \right\vert\geq 3$ and $ \left\vert G_2 \right\vert\geq 2 $. Suppose $G$ is generated by the union $G_1\cup G_2$. If there is a $G$-action on a set $X$ such that there are nonempty subsets $X_1,X_2\subseteq X$ with $X_2$ not contained in $X_1$ and such that for all $g\in G_i\setminus \set{e}$, we have $g.X_{j}\subseteq X_i$ when $i=1,2$ and $j=2,1$, then $G\cong G_1\ast G_2$.
\end{theorem}
\begin{proof}
  Suppose $w$ is a nontrivial reduced word in $G_1$ and $G_2$; by contracting if necessary, we start by assuming that $w$ is of the form $w = a_1b_1\cdots b_{k-1}a_k$ with $a_i\in G_1\setminus \set{e}$ and $b_i\in G_2\setminus \set{e}$. Our goal is to show that $w$ is not the identity in $g$. Observe that we have
  \begin{align*}
    w\left( X_2 \right) &= a_1b_1\cdots b_{k-1}a_k\left( X_2 \right)\\
                        &\subseteq a_1b_1\cdots b_{k-1}\left( X_1 \right)\\
                        &\vdots\\
                        &\subseteq a_1b_1\left( X_1 \right)\\
                        &\subseteq a_1\left( X_2 \right)\\
                        &\subseteq X_1.
  \end{align*}
  Since $X_2\nsubseteq X_1$, it follows that $w\neq e_G$. The other cases for the representation of the reduced word $w$ follow from conjugating with an element of $G_1$.
\end{proof}
\subsubsection{Applications}%
The first application we will consider is the matrix group $\SL_2\left( \Z \right)$. The matrices
\begin{align*}
  a &= \begin{pmatrix}1 & 2 \\ 0 & 1\end{pmatrix}\\
  b &= \begin{pmatrix}1 & 0 \\ 2 & 1\end{pmatrix}
\end{align*}
generate a subgroup of $\SL_2\left( \Z \right)$ that is isomorphic to the free group $F_2$. Considering the subsets of $\R^2$ given by
\begin{align*}
  A &= \set{ \begin{pmatrix}x\\y\end{pmatrix} : \left\vert x \right\vert > \left\vert y \right\vert }\\
  B &= \set{ \begin{pmatrix}x\\y\end{pmatrix} : \left\vert x \right\vert < \left\vert y \right\vert }.
\end{align*}
Then, we have
\begin{align*}
  a^{n} \begin{pmatrix}x\\y\end{pmatrix} &= \begin{pmatrix}x + 2ny\\y\end{pmatrix},
\end{align*}
whose action on $B$ satisfies
\begin{align*}
  \left\vert x + 2ny \right\vert &\geq \left\vert 2ny \right\vert - \left\vert x \right\vert\\
                                 &\geq 2\left\vert y \right\vert-\left\vert x \right\vert\\
                                 &> 2\left\vert y \right\vert-y\\
                                 &= \left\vert y \right\vert,
\end{align*}
so that $a^{n}.B \subseteq A$, and similarly $b^{m}.A \subseteq B$ for all $m,n\neq 0$. Thus, the subgroup $\left\langle a,b \right\rangle\subseteq \SL_2\left( \Z \right)$ is freely generated by the ping-pong lemma.
\begin{proposition}
  The subgroup $ \left\langle a,b \right\rangle\subseteq \SL_2\left( \Z \right) $ is finite index.
\end{proposition}
\section{Quasi-Isometry of Metric Spaces and Groups}%
\subsection{Quasi-Isometry of Metric Spaces}%
\begin{definition}
  Consider two metric spaces $\left( X,d \right)$ and $\left( Y,\rho \right)$, and a map $f\colon X\rightarrow Y$.
  \begin{itemize}
    \item We say $f$ is an isometric embedding if for all $x,x'\in X$, $\rho\left( f\left( x \right),f\left( x' \right) \right) = d\left( x,x' \right)$.
    \item We say $f$ is an isometry (or isometric isomorphism) if $f$ is an isometric embedding, and there exists an isometric embedding $g\colon Y\rightarrow X$ satisfying $f\circ g = \id_Y$ and $g\circ f = \id_X$.
    \item We say $f$ is a bilipschitz embedding if there is some constant $C > 0$ such that for all $x,x'\in X$,
      \begin{align*}
        \frac{1}{C}d\left( x,x' \right)\leq \rho\left( f\left( x \right),f\left( x' \right) \right)\leq Cd\left( x,x' \right).
      \end{align*}
    \item We say $f$ is a bilipschitz equivalence if it is a bilipschitz embedding and there exists a bilipschitz embedding $g\colon Y\rightarrow X$ satisfying $f\circ g = \id_Y$, $g\circ f = \id_X$. The spaces $X$ and $Y$ are then called bilipschitz equivalent.
    \item The map $f$ is called a $\left( C,K \right)$-quasi-isometric embedding if for all $x,x'\in X$, we have
      \begin{align*}
        \frac{1}{C}d\left( x,x' \right)-K \leq \rho\left( f\left( x \right),f\left( x' \right) \right)\leq Cd\left( x,x' \right) + K.
      \end{align*}
    \item The map $f'\colon X\rightarrow Y$ is said to have finite distance from $f$ if there is some $K > 0$ such that for all $x\in X$, we have $\rho\left( f(x),f'(x) \right) \leq K$.
    \item The map $f$ is called a quasi-isometry if there is a quasi-inverse quasi-isometric embedding $g\colon Y\rightarrow X$. That is, $g$ is a quasi-isometric embedding with $g\circ f$ having finite distance from $\id_X$ and $f\circ G$ having finite distance from $\id_Y$. We say $X$ and $Y$ are quasi-isometric, and write $X\sim_{\operatorname{QI}}Y$.
  \end{itemize}
\end{definition}
\begin{proposition}
  A map $f\colon (X,d)\rightarrow (Y,\rho)$ is a quasi-isometry if and only if it is a quasi-isometric embedding with quasi-dense image. A map is said to have quasi-dense image if there is some constant $C$ such that for all $y\in Y$, there is some $x\in X$ such that
  \begin{align*}
    \rho\left( f(x),y \right) &\leq C.
  \end{align*}
\end{proposition}
\begin{proof}
  If $f$ is a quasi-isometry, then there is a quasi-inverse quasi-isometric embedding $g\colon Y\rightarrow X$, so there is some constant $C$ such that for all $y\in Y$, $\rho\left( f\circ g(y),y \right) \leq c$; in particular, $f$ has quasi-dense image (choose $g(y)$ for the particular choice of $x$).

  Now, suppose $f$ is a quasi-isometric embedding with quasi-dense image. We will construct a quasi-inverse quasi-isometric embedding. We may find a constant $C$ such that for all $x,x'\in X$, we have
  \begin{align*}
    \frac{1}{C}d\left( x,x' \right)-C \leq \rho\left( f(x),f(x') \right)\leq Cd(x,x') + C
  \end{align*}
  and for all $y\in Y$, there exists $x\in X$ such that
  \begin{align*}
    \rho\left( f(x),y \right) &\leq C.
  \end{align*}
  For each $y$, we may then choose $x_y$ such that the map $g\colon Y\rightarrow X$ taking $y\mapsto x_y$ has $d\left( f\left( x_y \right),y \right)\leq C$ holding for all $y\in Y$.

  We claim that the map $g$ is quasi-inverse to $f$. We observe that by definition, we have
  \begin{align*}
    \rho\left( f\circ g(y),y \right) &= \rho\left( f\left( x_y \right),y \right)\\
                                     &\leq C.
  \end{align*}
  Conversely, using the fact that $f$ is a quasi-isometric embedding, we have for all $x\in X$ that
  \begin{align*}
    d\left( g\circ f(x),x \right) &= d\left( x_{f(x)},x \right)\\
                                  &\leq C\rho\left( f\left( x_{f(x)} \right),f(x) \right) + C^2\\
                                  &\leq 2C^2,
  \end{align*}
  meaning that both $f\circ g$ and $g\circ f$ have finite distance from the respective identity maps.

  Finally, we will show that $g$ is a quasi-isometric embedding. A slight modification of this argument will show that quasi-inverses of quasi-isometric embeddings are quasi-isometric embeddings.
  \begin{align*}
    d\left( g(y),g(y') \right) &= d\left( x_y,x_{y'} \right)\\
                               &\leq C\rho\left( f\left( x_y \right),f\left( x_{y'} \right) \right) + C^2\\
                               &\leq C\left( \rho\left( f\left( x_y \right),y \right) + \rho\left( y,y' \right) + \rho\left( f\left( x_{y'} \right),y' \right) \right) + C^2\\
                               &\leq C\left( \rho\left( y,y' \right) + 2C \right) + C^2\\
                               &= C\rho\left( y,y' \right) + 3C^2\\
    d\left( g\left( y \right),g\left( y' \right) \right) &= d\left( x_y,x_{y'} \right)\\
                                                         &\geq \frac{1}{C}\rho\left( f\left( x_y \right),f\left( x_{y'} \right) \right) - c^2\\
                                                         &\geq \frac{1}{C}\left( \rho\left( y,y' \right) - \rho\left( f\left( x_y \right),y \right) - \rho\left( f_{x_{y'}},y' \right) \right) - c^2\\
                                                         &\geq \frac{1}{C}\rho\left( y,y' \right) - \left(2+ c^2\right).
  \end{align*}
\end{proof}
The following inheritance properties can be computed directly from their respective definitions.
\begin{itemize}
  \item Every map at a finite distance from a quasi-isometric embedding (respectively, quasi-isometry) is a quasi-isometric embedding (respectively, quasi-isometry).
  \item If $f,f'\colon X\rightarrow Y$ are maps that are finite distance from each other, then if $g\colon Z\rightarrow X$ is any map, $f\circ g$ and $f'\circ g$ have finite distance from each other. If $g\colon Y\rightarrow Z$ is a quasi-isometric embedding, then $g\circ f$ and $g\circ f'$ also have finite distance from each other.
  \item Compositions of quasi-isometric (respectively, bilipschitz) embeddings are quasi-isometric (respectively, bilipschitz) embeddings.
  \item Compositions of quasi-isometries (respectively, bilipschitz equivalences) are quasi-isometries (respectively, bilipschitz equivalences).
\end{itemize}
Given a metric space $X$, the quasi-isometry group of $X$, denoted $ \operatorname{QI}(X) $, is the set of all quasi-isometries $f\colon X\rightarrow X$ modulo finite distance, equipped with the composition operation.
\subsection{Quasi-Isometry of Groups}%
\begin{definition}
  Let $X = \left( V,E \right)$ be a connected graph. The distance between $v,w\in V$ is equal to the length of the smallest path between $v$ and $w$. This is often known as the graph metric. Note that it is a metric on the vertex sets.
\end{definition}
It can be observed that an isometric isomorphism of vertex sets equipped with this metric is equivalent to an isometric isomorphism of the underlying graphs.
\begin{definition}
  Let $G$ be a finitely generated group with generating set $S$. The \textit{word length} of $g\in G$ with respect to $S$, denoted $\ell_S(g)$, is the smallest $n\geq 0$ such that $g$ may be written as a product of $n$ elements of $S$. That is,
  \begin{align*}
    \ell_S(g) &= \min\set{n : g = s_1s_2\cdots s_n\text{ with }s_i\in S\text{ and }n\geq 0}.
  \end{align*}
  The \textit{word metric} is then
  \begin{align*}
    d_S(g,h) &= \ell_S\left( g^{-1}h \right).
  \end{align*}
\end{definition}
\begin{proposition}
  The word metric is a left-invariant metric. That is, it satisfies the following:
  \begin{enumerate}[(i)]
    \item $d_S\left( g,h \right) = 0$ if and only if $g = h$;
    \item $d_S\left( g,h \right) = d_S\left( h,g \right)$;
    \item $d_S\left( g,h \right) \leq d_S\left( g,k \right) + d_S\left( k,h \right)$;
    \item $d_S\left( g,h \right) = d_S\left( kg,kh \right)$.
  \end{enumerate}
\end{proposition}
\begin{proof}\hfill
  \begin{enumerate}[(i)]
    \item We have that $\ell\left( g \right) = 0$ if and only if $g = e$ by the fact that any nontrivial element must be a nontrivial product of elements of $S$.
    \item Notice that $\ell\left( g \right) = \ell\left( g^{-1} \right)$ since $S$ is symmetric, so if
      \begin{align*}
        g &= s_1\cdots s_n
          \intertext{realizes $\ell(g)$, then}
        g^{-1} &= s_n^{-1}\cdots s_1^{-1}
      \end{align*}
      must realize $\ell\left( g^{-1} \right)$, as otherwise we would have a smaller representation for $g$ in terms of a product of elements of $S$.
    \item Observe that, by the existence of relations and word reductions, we have $\ell\left( gh \right)\leq \ell\left( g \right) + \ell\left( h \right)$.
    \item We have
      \begin{align*}
        d\left( kg,kh \right) &= \ell\left( \left( kg \right)^{-1}\left( kh \right) \right)\\
                              &= \ell\left( g^{-1}k^{-1}kh \right)\\
                              &= \ell\left( g^{-1}h \right)\\
                              &= d\left( g,h \right).
      \end{align*}
  \end{enumerate}
\end{proof}
\begin{proposition}
  Let $G$ be a finitely generated group, and let $S$ and $S'$ be finite generating sets of $G$. Then, the identity map is a bilipschitz equivalence of the metric spaces $\left( G,d_S \right)$ and $\left( G,d_{S'} \right)$.
\end{proposition}
\begin{proof}
  Since $S$ is finite, the value
  \begin{align*}
    c &= \max\set{s\in S} d_{S'} \left( e,s \right)
  \end{align*}
  is also finite. Suppose $g,h\in G$, and let $n = d_S\left( g,h \right)$. Then, $g^{-1}h = s_1\cdots s_n$ for some $s_1,\dots,s_n\in S$. By the triangle inequality and left-invariance of the metric $d_{S'}$, we obtain
  \begin{align*}
    d_{S'}\left( g,h \right) &= d_{S'}\left( g,gs_1\cdots s_n \right)\\
                             &\leq d_{S'}\left( g,gs_1 \right) + \cdots + d_{S'} \left( gs_1\cdots s_{n-1},gs_1\cdots s_n \right)\\
                             &= d_{S'}\left( e,s_1 \right) + d_{S'}\left( e,s_2 \right) + \cdots + d_{S'}\left( e,s_n \right)\\
                             &\leq cn\\
                             &= cd_S\left( g,h \right).
  \end{align*}
  Symmetry gives the desired bilipschitz equivalence.
\end{proof}
Due to the word metric, we are able to define notions of bilipschitz equivalence and quasi-isometry for finitely generated groups.
\begin{definition}
  Let $G$ be a finitely generated group.
  \begin{enumerate}[(i)]
    \item We say $G$ is bilipschitz equivalent to a metric space $X$ if for some finite generating set $S$, the metric spaces $\left( G,d_S \right)$ and $X$ are bilipschitz equivalent.
    \item We say $G$ is quasi-isometric to a metric space $X$ if for some finite generating set $S$, the metric spaces $\left( G,d_S \right)$ and $X$ are quasi-isometric. We will write $G\sim_{\operatorname{QI}}X$ if $G$ and $X$ are quasi-isometric.
  \end{enumerate}
\end{definition}
\begin{remark}\hfill
  \begin{enumerate}[(i)]
    \item Bijective quasi-isometries of finitely generated groups are bilipschitz equivalences due to the fact that the minimum distance between any two group elements is $1$.
    \item If $S\subseteq G$ is a generating set for some group $G$, then $S$ is finite if and only if the word metric $d_S$ is proper, in the sense that the balls of finite radius in $\left( G,d_S \right)$ are finite.
  \end{enumerate}
\end{remark}
We start by considering the quasi-isometry classification of finite groups. It turns out that a finitely generated group is quasi-isometric to a finite group if and only if it is finite. This follows from the fact that all finite groups yield metric spaces of finite diameter, and all finite diameter metric spaces are quasi-isometric to each other; conversely, a group is quasi-isometric to a finite group if and only if there is a finite generating set that yields a metric space with finite diameter, but since balls are finite, it follows that the group is itself finite.

Next, we claim that the groups $\Z$, $\Z\times \Z/2\Z$, and $D_{\infty}\cong \Z\rtimes \Z/2\Z$ are bilipschitz equivalent. Consider the equivalent presentations
\begin{align*}
  D_{\infty} &= \left\langle x,y | x^2,y^2\right\rangle\\
             &= \left\langle s,t | t^2,tst^{-1}=s^{-1} \right\rangle.
\end{align*}
Then, the Cayley graph $ \operatorname{Cay}\left( D_{\infty},\set{x,y} \right) $ is isomorphic to the Cayley graph $\operatorname{Cay}\left( \Z,\set{1} \right)$ --- map $\set{x,xy,xyx,\dots}$ to strictly positive numbers and map $\set{y,yx,yxy,\dots}$ to strictly negative numbers.

On the other hand, the Cayley graph $\operatorname{Cay}\left( D_{\infty},\set{s,t} \right)$ is isomorphic to the Cayley graph $\operatorname{Cay}\left( \Z\times \Z/2\Z,\set{\left( 1,[0] \right),\left( 0,[1] \right)} \right)$. Since the word metrics on $D_{\infty}$ with respect to $\set{x,y}$ and $\set{s,t}$ are bilipschitz equivalent, we get our desired result.

In fact, it is possible to show something a lot more rigid.
\begin{theorem}
  A finitely generated group is quasi-isometric to $\Z$ if and only if there is a finite-index subgroup isomorphic to $\Z$.
\end{theorem}
\subsection{Quasi-Geodesics}%
\begin{definition}
  Let $\left( X,d \right)$ be a metric space.
  \begin{enumerate}[(i)]
    \item A geodesic of length $L\geq 0$ is an isometric embedding $\gamma\colon [0,L]\rightarrow X$, where $[0,L]$ is endowed with the standard metric on $\R$. The point $\gamma(0)$ is the start point of $\gamma$ and the point $\gamma(L)$ is the endpoint of $\gamma$.
    \item The metric space $X$ is called geodesic if for all $x,x'\in X$ there exists a geodesic in $X$ with start point $x$ and end point $x'$.
    \item Given $c > 0$ and $b\geq 0$, a $\left( c,b \right)$-quasi-geodesic is a $\left( c,b \right)$-quasi-isometric embedding $\gamma\colon [t,t']\rightarrow X$. The point $\gamma(t)$ is the start point of $\gamma$ and the point $\gamma(t')$ is the end point of $\gamma$.
    \item The space $\left( X,d \right)$ is called $\left( c,b \right)$-quasi-geodesic if for all $x,x'\in X$, there exists a $\left( c,b \right)$-quasi-geodesic with start point $x$ and end point $x'$.

      More concretely, for any $x,x'\in X$, there is a map $\gamma\colon [t,t']\rightarrow X$ such that for any $s,s'\in [t,t']$, we have
      \begin{align*}
        \frac{1}{c}\left\vert s-s' \right\vert - b \leq d\left( \gamma(s),\gamma(s') \right) \leq c\left\vert s-s' \right\vert + b.
      \end{align*}
  \end{enumerate}
\end{definition}
The most important example of a quasi-geodesic space is that of a connected graph. Specifically, if $X = \left( V,E \right)$ is a connected graph, then the graph metric on $V$ (defined as the shortest path between two vertices) turns $V$ into a $\left( 1,1 \right)$-quasi-geodesic space.

It is possible to turn a graph into a geodesic space by gluing any two vertices connected by an edge using a unit interval. In fact, this applies to any quasi-geodesic metric space, and the resulting geodesic space is quasi-isometric to the original quasi-geodesic space.
\subsection{Milnor--Schwarz Lemma}%
We wish to understand what finitely-generated groups look like up to quasi-isometry. The Milnor--Schwarz lemma helps establish this via a sufficiently nice action of the group.
\begin{proposition}[Milnor--Schwarz Lemma]
  Let $G$ be a group, and let $G$ act on a nonempty metric space $\left( X,d \right)$ by isometries.

  Suppose there are constants $c,b > 0$ such that $X$ is $\left( c,b \right)$-quasi-geodesic, and suppose there is a subset $B\subseteq X$ with:
  \begin{itemize}
    \item $\operatorname{diam}\left( B \right)$ finite;
    \item $\bigcup_{g\in G}g.B = X$;
    \item the action of $G$ satisfies $ \left\vert \set{g\in G : g.B'\cap B' \neq \emptyset} \right\vert $ is finite, where we define
      \begin{align*}
        B' &\coloneq \set{x\in X : \text{ there exists $y\in B$ such that }d\left( x,y \right)\leq 2b}.
      \end{align*}
  \end{itemize}
  Then, the following hold.
  \begin{enumerate}[(i)]
    \item The group $G$ is generated by $S$, and in particular, $G$ is finitely generated.
    \item For all $x\in X$, the map $g\mapsto g.x$ yields a quasi-isometry with respect to the word metric $d_S$ on $G$.
  \end{enumerate}
\end{proposition}
\begin{proof}
  We start by showing that $S$ generates $G$. Let $g\in G$, and let $x\in B$. Since $X$ is $\left( c,b \right)$-quasi-geodesic, there is a $\left( c,b \right)$-quasi-geodesic map $\gamma\colon [0,L]\rightarrow X$ from $x$ to $g.x$.

  Let $n = \lceil Lc/b\rceil$. For each $j\in \set{0,\dots,n-1}$, we define $t_j = j\frac{b}{c}$, $t_n = L$, and $x_j = \gamma\left( t_j \right)$. Notice that $x_0 = x$ and $x_n = g.x$. Since the translates of $B$ cover all of $X$, there are group elements $g_j\in G$ with $x_j\in g_j.B$. We may choose $g_0 = e$ and $g_n = g$.

  Then, for each $j=1,\dots,n$, the element $s_j = g_{j-1}^{-1}g_j$ lies in $S$. Since $\gamma$ is a $\left( c,b \right)$-quasi-geodesic, it follows that
  \begin{align*}
    d\left( x_{j-1},x_j \right) &\leq c\left\vert t_{j-1}-t_j \right\vert + b\\
                                &\leq c \frac{b}{c} + b\\
                                &\leq 2b.
  \end{align*}
  In particular, $x_j\in g_{j-1}.B'$, while the definition of $x_j$ has $x_j\in g_j.B \subseteq g_j.B'$, whence $g_{j-1}.B'\cap g_j.B' \neq \emptyset$. That is, it follows that $s_{j} = g_{j-1}^{-1}g_j\in S$.

  Therefore, we have
  \begin{align*}
    g &= g_n\\
      &= g_{n-1}g_{n-1}^{-1}g_n\\
      &= g_{n-1}s_{n-1}g_n\\
      &\vdots\\
      &= s_1\cdots s_n
  \end{align*}
  is contained in the subgroup generated by $S$.

  Now, we will show that $G$ is quasi-isometric to $X$. We will show that $g\xmapsto{\varphi} g.x$ is a quasi-isometric embedding with quasi-dense image. Since $G$ acts by isometries on $X$ and the $G$-translates of $B$ cover all of $X$, it follows that we may assume that $x\in B$.

  First, we will show that $g\mapsto g.x$ has quasi-dense image. Notice that if $x'\in X$, then there is some $g\in G$ with $x'\in g.B$, so since $g.x\in g.B$, we have
  \begin{align*}
    d\left( x',\varphi(g) \right) &= d\left( x',g.x \right)\\
                                  &\leq \operatorname{diam}\left( g.B \right)\\
                                  &= \operatorname{diam}\left( B \right).
  \end{align*}
  Since we assumed that $ \operatorname{diam}\left( B \right) $ is finite, this yields a quasi-dense image.

  Now, we will show that $\varphi$ is a quasi-isometric embedding. Let $g\in G$ be given by $g = s_1\cdots s_n$ realizing the word length of $g$ (i.e., the distance to the identity). We will establish a uniform lower bound on $d\left( \varphi(e),\varphi(g) \right)$ in terms of $d_S\left( e,g \right)$. Let $\gamma\colon [0,L]\to X$ be a $\left( c,b \right)$-quasi-geodesic from $x$ to $g.x$. Then, we have
  \begin{align*}
    d\left( \varphi(e),\varphi(g) \right) &= d\left( x,g.x \right)\\
                                          &= d\left( \gamma(0),\gamma(L) \right)\\
                                          &\geq \frac{1}{c}L-b\\
                                          &\geq \frac{1}{c}\frac{b\left( n-1 \right)}{c} - b\\
                                          &= \frac{b}{c^2}n - \frac{b}{c^2}b\\
                                          &\geq \frac{b}{c^2}d_S\left( e,g \right) - \frac{b}{c^2}-b.
  \end{align*}
  Conversely, we use the triangle inequality, the fact that $G$ acts via isometries, and that $s_j.B'\cap B' \neq \emptyset$ to obtain
  \begin{align*}
    d\left( \varphi(e),\varphi(g) \right) &= d\left( x,g.x \right)\\
                                          &\leq d\left( x,s_1.x \right) + \cdots + d\left( \left( s_1\cdots s_{n-1} \right).x, g.x \right)\\
                                          &= d\left( x,s_1.x \right) + \cdots + d\left( x,s_n.x \right)\\
                                          &\leq 2n\left( \operatorname{diam}\left( B \right) + 2b \right)\\
                                          &= 2\left( \operatorname{diam}\left( B \right) + 2b \right)d_S\left( e,g \right).
  \end{align*}
  Notice that $d\left( \varphi(g),\varphi(h) \right) = d\left( \varphi(e),\varphi\left( g^{-1}h \right) \right)$ and $d_S\left( g,h \right) = d_S\left( e,g^{-1}h \right)$ holds for all $g,h\in G$, so that $\varphi$ is thus a quasi-isometric embedding.
\end{proof}
We next establish a more topological (and useful) formulation of the Milnor--Schwarz Lemma.
\begin{corollary}[Topological Formulation of Milnor--Schwarz]
  Let $G$ be a group acting by isometries on a nonempty proper geodesic metric space $\left( X,d \right)$. Suppose this action is proper and cocompact. Then, $G$ is finitely generated, and for all $x\in X$, the map $g\mapsto g.x$ is a quasi-isometry.
\end{corollary}
First, we recall/define the most important ideas from this formulation.
\begin{itemize}
  \item A metric space $X$ is called proper if for all $x\in X$ and all $r > 0$, the closed ball $B\left( x,r \right)$ is compact with respect to the metric topology.
  \item An action $G\curvearrowright X$ is called proper if for all compact $B\subseteq X$, the set $\set{g\in G : g.B\cap B \neq \emptyset}$ is finite.
  \item An action $G\curvearrowright X$ is called cocompact if the (left) quotient space $G\backslash X$ is compact with respect to the quotient topology.
\end{itemize}
\begin{example}\hfill
  \begin{itemize}
    \item The translation action of $\Z$ on $\R$ is proper. Furthermore, this action is cocompact since the quotient is given by $S^1$
    \item The action by deck transformations of the fundamental group of a locally compact path-connected topological space on the universal cover is proper. If $X$ is compact, then the action by deck transformations is also cocompact since the quotient is given by $ \pi_1\left( X,x_0 \right)\backslash \widetilde{X}\cong X $.
    \item Stabilizer groups of a proper action are finite, but the converse is not necessarily true. For instance, the action of $\Z$ on $S^1$ by irrational rotation is free but not proper.
    \item The translation action of $\Z$ on $\R^2$ is not cocompact since the quotient is given by $S^1\times \R$.
    \item The action of $\SL_2\left( \Z \right)$ by Möbius transformations on the upper half-plane $ \mathbb{H} $ is not cocompact.
  \end{itemize}
\end{example}
\begin{proof}[Proof of Topological Formulation]
  Since $X$ is geodesic, it follows that $X$ is $\left( 1,b \right)$-quasi-geodesic for all $b > 0$. Our goal is to find a suitable subset $B\subseteq X$ that allows us to apply the original Milnor--Schwarz Lemma.

  Since the projection $\pi\colon X\rightarrow G\backslash X$ is an open map, and $G\backslash X$ is compact, it follows that there is a subspace $B\subseteq X$ such that $\pi(B) = G\backslash X$. Thus, we have $\bigcup_{g\in G}g.B = X$, and 
  \begin{align*}
    B' &= \set{y\in X : d\left( x,y \right) \leq 2b\text{ for some }x\in B}
  \end{align*}
  has finite diameter. Since $X$ is a proper space, the subset $B'$ is compact. Since the action of $G$ on $X$ is proper, it follows that $\set{g\in G : g.B'\cap B' \neq \emptyset}$ is finite. Thus, the conditions for the Milnor--Schwarz lemma apply.
\end{proof}
\subsection{Applications of the Milnor--Schwarz Lemma}%
\begin{corollary}
  Finite index subgroups of finitely generated groups are finitely generated and quasi-isometric to the ambient group.
\end{corollary}
\begin{proof}
  Let $H\subseteq G$ be a finite index subgroup of the finitely generated group $G$. The left translation action of $H$ on $\left( G,d_S \right)$ for some finite generating set $S$ satisfies the conditions of the Milnor--Schwarz lemma.
  \begin{itemize}
    \item The action of $H$ on $\left( G,d_S \right)$ is isometric.
    \item The space $\left( G,d_S \right)$ is $\left( 1,1 \right)$-quasi-geodesic.
    \item If $B\subseteq G$ is a finite set of representatives for the right cosets $H\backslash G$, then $HB = G$, and
      \begin{align*}
        B' &= \set{y\in G : d_S\left( x,y \right)\leq 2\text{ for some }x\in B}
      \end{align*}
      is finite, and the set
      \begin{align*}
        T &= \set{h\in H : h.B'\cap B'\neq \emptyset}
      \end{align*}
      is also finite.
  \end{itemize}
  Thus, $H$ is finitely generated by $T$ and the inclusion $H\hookrightarrow G$ is a quasi-isometry with respect to any word metric.
\end{proof}
This yields a very significant notion of equivalence between groups.
\begin{definition}\hfill
  \begin{itemize}
    \item Two groups $G$ and $H$ are called commensurable if they contain finite index subgroups $G'\subseteq G$ and $H'\subseteq H$ with $G'\cong H'$.
    \item Two subgroups $G$ and $H$ are weakly commensurable if they contain finite index subgroups $G'\subseteq G$ and $H'\subseteq H$ such that there exist finite normal subgroups $N\trianglelefteq G'$ and $M\trianglelefteq H'$ with isomorphic quotients.
  \end{itemize}
\end{definition}
\begin{corollary}
  Let $G$ be a group.
  \begin{enumerate}[(i)]
    \item If $G'\subseteq G$ is a finite index subgroup of $G$, then $G'$ is finitely generated if and only if $G$ is finitely generated.
    \item Let $N$ be a finite normal subgroup. Then, $G/N$ is finitely generated if and only if $G$ is finitely generated, and $G\sim_{\operatorname{QI}} G/N$.
  \end{enumerate}
\end{corollary}
\begin{proof}\hfill
  \begin{enumerate}[(i)]
    \item Using the first corollary to Milnor--Schwarz, it suffices to show that $G$ is finitely generated if $G'$ is. Yet, since $G'\subseteq G$ is finite index, it follows that taking the representatives of the finitely many cosets as well as the generators of $G'$ yields a generating set for $G$.
    \item If $G$ is finitely generated, then so is $G/N$. Conversely, if $G/N$ is finitely generated, then taking representatives $s'\in \pi^{-1}\left( sN \right)$ for each $sN$ in the generating set for $G/N$, along with the finite set $N$, yields a generating set for $G$.

      If $G$ and $G/N$ are finitely generated, and $S$ is a finite generating set for $G/N$, then on $\left( G/N,d_S \right)$, we may define an isometric action of $G$ by taking $g.x\coloneq \pi(g)x$ for any $x\in G/N$. This action satisfies the conditions of the Milnor--Schwarz lemma, so that $G\sim_{\operatorname{QI}}G/N$.
  \end{enumerate}
\end{proof}
In particular, this means that if two groups are weakly commensurable, then they're quasi-isometric. However, this doesn't go in reverse.

For instance, if $F_3$ denotes the free group of rank $3$ and $F_4$ the free group of rank $4$, then it can be shown that the finitely generated groups $\left( F_3\times F_3 \right)\ast F_3$ and $\left( F_3\times F_3 \right)\ast F_4$ are bilipschitz equivalent, but they are not commensurable since the Euler characteristic of the former is zero while the Euler characteristic of the latter is nonzero.

Now, we will discuss some applications of the Milnor--Schwarz lemma to Riemannian geometry.
\begin{corollary}
  Let $M$ be a closed connected Riemannian manifold with Riemannian universal cover $ \widetilde{M} $. Then, $\pi_1(M)$ is finitely generated and the action of $\pi_1(M)$ on $ \widetilde{M} $ by deck transformations yields a quasi-isometry between $\pi_1(M)$ and $ \widetilde{M} $.
\end{corollary}
\begin{proof}[Proof Sketch]
  First, one uses arguments from Riemannian geometry to show that $ \widetilde{M} $ is a proper geodesic metric space, and the action of $ \pi_1(M) $ on $ \widetilde{M} $ is isometric, proper (fundamental groups act properly on the universal cover by deck transformations), and cocompact (since the quotient is given by the compact space $M$). Therefore, the topological formulation of the Milnor--Schwarz lemma yields the final proof.
\end{proof}
\begin{definition}
  We call a Riemannian manifold flat if its Riemannian universal cover is isometrically isomorphic to the Euclidean space $\R^n$ of the same dimension.

  We call a Riemannian manifold hyperbolic if its Riemannian universal cover is isometrically isomorphic to the hyperbolic space $ \mathbb{H}^n $ of the same dimension.
\end{definition}
These definitions allow us to yield the following classification.
\begin{corollary}\hfill
  \begin{enumerate}[(i)]
    \item If $M$ is a closed connected flat Riemannian $n$-manifold, then its fundamental group is quasi-isometric to $\R^n$, hence quasi-isometric to $\Z^n$.
    \item If the fundamental group of a closed connected smooth $n$-manifold is not quasi-isometric to $\R^n$ or $\Z^n$, then the manifold does not admit a flat Riemannian metric.
    \item If $M$ is a closed connected hyperbolic Riemannian manifold, then $\pi_1(M)$ is quasi-isometric to the hyperbolic space $ \mathbb{H}^n $.
    \item If the fundamental group of a closed connected smooth $n$-manifold is not quasi-isometric to $ \mathbb{H}^n $, then this manifold does not admit a hyperbolic Riemannian metric.
  \end{enumerate}
\end{corollary}
Finally, we discuss the dynamic criterion for quasi-isometry.
\begin{definition}
  Let $G$ and $H$ be groups. A set-theoretic coupling for $G$ and $H$ is a nonempty set $X$ with a left action $G\curvearrowright X$ and right action $X\curvearrowleft H$ that commute with each other, in that $\left( g.x \right).h = g.\left( x.h \right)$, such that $X$ admits a subset $K$ satisfying the following.
  \begin{itemize}
    \item The $G$ and $H$ translates of $K$ cover $K$. That is, $G.K = K.H$.
    \item The sets
      \begin{align*}
        F_G &= \set{g\in G : g.K\cap K \neq \emptyset}
        F_H &= \set{h\in H : K.h\cap H \neq \emptyset}
      \end{align*}
      are finite.
    \item For each $g\in G$, there is a finite subset $F_H(g)\subseteq H$ with $g.K\subseteq K.F_H(g)$, and for each $h\in H$, there is a finite subset $F_G(h)\subseteq H$ such that $K.h\subseteq F_G(h).K$.
  \end{itemize}
\end{definition}
For instance, if $G,H\subseteq X$ are finite-index subgroups acting by multiplication on the left and right respectively, then $G$ and $H$ admit a set-theoretic coupling since multiplication in $X$ is associative and the subset $K$ can be given by taking the finite sets of representatives for the right cosets of $G$ and left cosets of $H$.\footnote{The $G$-orbit of an element $x\in X$ is given by $Gx$ and the $H$-orbit of $x$ is given by $xH$.}
\begin{definition}
  Let $G$ and $H$ be groups. A topological coupling for $G$ and $H$ is a nonempty locally compact space $X$ with a proper cocompact left action $G\curvearrowright X$ by homeomorphisms and a proper cocompact right action $X\curvearrowleft H$ by homeomorphisms that commute with each other.
\end{definition}
\begin{theorem}[Gromov's Coupling Criterion for Quasi-Isometry]
  Let $G$ and $H$ be finitely generated groups. The following are equivalent:
  \begin{enumerate}[(i)]
    \item there is a topological coupling for $G$ and $H$;
    \item there is a set-theoretic coupling for $G$ and $H$;
    \item the groups $G$ and $H$ are quasi-isometric.
  \end{enumerate}
\end{theorem}
\begin{proof}
  First, assuming (i), we observe that there is a compact subset $K\subseteq X$ such that $G.K = X = K.H$. Since the actions of $G$ and $H$ are proper, the sets
  \begin{align*}
    F_G &= \set{g\in G : g.K\cap K \neq \emptyset}\\
    F_H &= \set{h\in H : K.h\cap K\neq \emptyset}
  \end{align*}
  are finite. Compactness of $K$ and the local compactness of $X$ yield that for every $g\in G$, there is a finite subset $F_H(g)\subseteq H$ such that $g.K\subseteq K.F_H(g)$, and similarly for elements of $H$. Thus, we obtain a set-theoretic coupling.

  Now, we will show that if $G$ and $H$ admit a set-theoretic coupling, $G$ and $H$ are quasi-isometric. Given $x\in K\subseteq X$, we obtain a map $f\colon G\rightarrow H$ such that $g^{-1}.x \in K.f(g)^{-1}$ for all $g\in G$. Given finite generating sets $S\subseteq G$ and $T\subseteq H$, then for any subsets $A\subseteq G$ and $B\subseteq H$, we define
  \begin{align*}
    D_TB &\coloneq \sup_{b\in B} d_T(e,b)\\
    D_SA &\coloneq \sup_{a\in A} d_S(e,a).
  \end{align*}

\end{proof}
\section{Amenability of Groups}%
We are only concerned with discrete groups here. 
\begin{definition}
  A group $G$ is called \textit{amenable} if there exists a map $m\colon \ell^{\infty}\left( G \right)\rightarrow \C$ (known as a mean) satisfying
  \begin{itemize}
    \item $m(1) = 1$ (that is, the map $1\colon G\rightarrow \C$ taking all elements to $1$);
    \item $m\left( f \right) \geq 0$ for all positive $f\in \ell^{\infty}\left( G \right)$;
    \item for all $g\in G$ and all $f\in \ell^{\infty}\left( G \right)$, we have $m\left( \lambda_g(f) \right) = m(f)$, where
      \begin{align*}
        \lambda_g(f)(s) &= f\left( g^{-1}s \right)
      \end{align*}
      is the left regular action of $G$ on itself.
  \end{itemize}
\end{definition}
\begin{proposition}
  Finite groups are amenable.
\end{proposition}
\begin{proof}
  The averaging operator
  \begin{align*}
    f\mapsto \frac{1}{\left\vert G \right\vert} \sum_{g\in G} f\left( g \right)
  \end{align*}
  is a mean.
\end{proof}
To show that abelian groups are amenable, we recall the Markov--Kakutani Fixed-Point Theorem.
% Coamenability

\section{Four Perspectives on Soficity}%
We will work to understand soficity (a type of generalized approximation property for groups) with a variety of perspectives.
\subsection{The Law of the Group von Neumann Algebra}%
This is the approach that my advisor, Ben Hayes, takes towards tracial von Neumann algebras, specialized to the case of groups.

A tracial von Neumann algebra is a pair $\left( M,\tau \right)$, where $M$ is a von Neumann algebra and $\tau\colon M\rightarrow \C$ is a linear functional that satisfies:
\begin{itemize}
  \item positivity: $\tau\left( x^{\ast}x \right) \geq 0$ for all $x\in M$;
  \item normalization: $\tau\left( 1 \right) = 1$;
  \item traciality: for all $x,y\in M$, $\tau\left( xy \right) = \tau\left( yx \right)$.
\end{itemize}
If $I$ is any index set, an $I$-tuple $a\in M^{I}$ yields a $\ast$-algebra homomorphism from the space of all noncommutative $\ast$-polynomials, denoted $ \C^{\ast}\left\langle \left( T_i \right)_{i\in I} \right\rangle $, given by $P(a)\in M$ for all $P\in \C^{\ast}\left\langle \left( T_i \right)_{i\in I} \right\rangle$.

The \textit{law} of an $I$-tuple $a\in M^{I}$ is then a linear functional $\ell_a\colon \C^{\ast}\left\langle \left( T_i \right)_{i\in I} \right\rangle\rightarrow \C$ given by
\begin{align*}
  \ell_a\left( P \right) &= \tau\left( P(a) \right).
\end{align*}
Notice that in the case of a single-element set for $I$ and a normal or self-adjoint element for $A$, the law is given by integration with respect to the spectral measure of $a$,
\begin{align*}
  \ell_a(P) &= \int_{}^{} P\:d\mu_a.
\end{align*}
In a sense, the law is a generalization of the spectral measure.

Now, consider a group $G$ with finite symmetric generating set $S$. Then, in the group von Neumann algebra $L(G)$, there is a distinguished tuple $\left( \lambda_s \right)_{s\in S}\in L(G)^{S}$ given by the left-translation operators $\lambda_s$ for each $s\in S$. The corresponding law is then given by
\begin{align*}
  \ell_S\left( P \right) &= \iprod{P\left( \left( \lambda_s \right)_{s\in S} \right)\delta_e}{\delta_e}.
\end{align*}
We start by viewing the symmetric group $\Sym(n)\subseteq U\left( \M_n\left( \C \right) \right)$ as permutation matrices. Then, any $S$-tuple $\sigma\in \sym(n)^{S}$ has law given by
\begin{align*}
  \ell_{\sigma}\left( P \right) &= \tr_n\left( P\left( \left( \sigma_s \right)_{s\in S} \right) \right).
\end{align*}
Notice that, given a word $w$ in $\F_S$, which we can view as a monomial in $\C^{\ast}\left\langle \left( T_i \right)_{i\in I} \right\rangle$, we observe that the non-normalized trace is the cardinality of the fixed-point set, so that
\begin{align*}
  \tr_n\left( w\left( \left( \sigma_s \right)_{s\in S} \right) \right) &= \frac{1}{n}\left\vert \set{j\in [n] : w\left( \left( \sigma_{s} \right)_{s\in S} \right)(j) = j} \right\vert.
\end{align*}
We then say that $G$ is sofic if and only if there is a sequence of natural numbers $\left( k_n \right)_{n}$ and a corresponding sequence of $S$-tuples $\left( \sigma_n \right)_{n}$ with each $\sigma_n\in \sym\left( k_n \right)^{S}$, such that the laws $\ell_{\sigma_n}$ converge weak* (i.e., pointwise) to the law $\ell_S$.

\subsection{Metric Ultraproducts}%
The original notion of an ultraproduct came from consideration of normed spaces and a development of nonstandard analysis.

Let $E_{\alpha}$ be a family of normed spaces indexed by a set $A$. We may form the $\ell^{\infty}$-direct product of the set $\set{E_{\alpha}\colon \alpha\in A}$ by taking
\begin{align*}
  \prod_{\alpha\in A} E_{\alpha} &\coloneq \set{\left( x_{\alpha} \right)_{\alpha\in A} : x_{\alpha}\in E_{\alpha}\text{ and }\sup_{\alpha}\norm{x_{\alpha}} < \infty},
\end{align*}
and we may equip it with a norm
\begin{align*}
  \norm{x} &= \sup_{\alpha\in A} \norm{x_{\alpha}}_{\alpha}.
\end{align*}
Write $ \mathcal{E} $ for the $\ell^{\infty}$-direct product.

Given a (non-principal) ultrafilter $ \mathcal{U} $ on $A$, there is a way to define a ``null subspace'' by taking
\begin{align*}
  \mathcal{N}_{\mathcal{U}} &= \set{x\in \mathcal{E} : \lim_{ \mathcal{U} }\norm{x_{\alpha}} = 0},
\end{align*}
where the limit along the ultrafilter $\mathcal{U}$ of a sequence $\left( y_{\alpha} \right)_{\alpha}$ is defined to be $y$ if for all $\ve > 0$, the set $\set{\alpha : \norm{y_{\alpha}-y} < \ve}$ is contained in $ \mathcal{U} $. The subspace $ \mathcal{N}_{\mathcal{U}} $  is closed. The metric ultraproduct of the family $\left( E_{\alpha} \right)_{\alpha\in A}$ is then given by
\begin{align*}
  \prod_{\mathcal{U}} E_{\alpha} &\coloneq \mathcal{E}/\mathcal{N}_{\mathcal{U}},
\end{align*}
with norm given by
\begin{align*}
  \norm{\left[ x \right]_{\mathcal{U}}} &= \lim_{\mathcal{U}} \norm{x_{\alpha}}.
\end{align*}
We write $E$ for the ultraproduct. When the ultrafilter is nonprincipal, then the ultraproduct is complete. If $\left( x_k \right)_k$ is a Cauchy sequence of elements of $ E $, then for each $i\in \N$, we may find some $N_i$ such that
\begin{align*}
  \norm{x_N' - x_N} &< 2^{-i}
\end{align*}
for all $N,N' \geq N_i$. For each $k$, we select a representative $\left( x_k(\alpha) \right)_{\alpha\in A}$, and define
\begin{align*}
  I_i &= \set{\alpha\in A : \norm{x_{N_i}(\alpha) - x_{N_{i+1}}(\alpha)} < 2^{-i}}.
\end{align*}
We have that each $I_i$ is contained in $ \mathcal{U} $, and without loss of generality we may assume that $I_1\supseteq I_2\supseteq\cdots$ with empty intersection, since the ultrafilter $ \mathcal{U} $ is non-principal, and thus is countably incomplete. Define an element $x\in \mathcal{E}$ by taking
\begin{align*}
  x|_{I_i\setminus I_{i+1}} &= x_{N_i}|_{I_i\setminus I_{i+1}}.
\end{align*}
The equivalence class $\left[ x \right]_{\mathcal{U}}$ is the limit of the Cauchy sequence $\left( x_k \right)_k$.

Next, we concern ourselves with ultraproducts of metric groups. Given a family $\left( G_{\alpha},d_{\alpha} \right)$ of topological groups with left-invariant metrics,\footnote{Such invariant metrics always exist for metrizable groups by a theorem of Kakutani.} we may form an ultraproduct by following an analogous process for the case of metric spaces, but the result is only a metric space on which the Cartesian product group $\prod_{\alpha\in A}G_{\alpha}$ acts by isometries (i.e., a homogeneous space), rather than a proper group.
\begin{example}
  Consider the symmetric group $\sym\left(\N\right)$ on the natural numbers. The standard Polish topology on $\sym(\N)$ is the subspace topology induced by the product topology on the prodiscrete space $\N^{\N}$. This is a separable group topology with a compatible metric given by
  \begin{align*}
    d\left( \sigma,\tau \right) &= \sum_{i=1}^{\infty} \set{2^{-i} : \sigma(i)\neq \tau(i)}.
  \end{align*}
  We note that this is a bounded analogue of the Hamming distance for finite symmetric groups.

  We will form an ultrapower of the group $\left( \sym(\N),d \right)$ with regard to a nonprincipal ultrafilter $ \mathcal{U} $ on $\N$. Since the metric $d$ is itself bounded, we see that the space $\mathcal{G} = \left( \sym(\N) \right)^{\N}$ is the analogue of the space $\mathcal{E}$ from the construction of the ultraproduct of normed spaces.

  The set
  \begin{align*}
    \mathcal{N} &= \set{x : \lim_{\mathcal{U}}d\left( x_{\alpha},e \right) = 0}
  \end{align*}
  is a subgroup of $ \mathcal{G} $, seeing as the invariance of the metric $d$ yields
  \begin{align*}
    d\left( xy, e\right) &= d\left( y,x^{-1} \right)\\
                         &\leq d\left( y,e \right) + d\left( x^{-1},e \right)\\
                         &= d\left( y,e \right) + d\left( x,e \right).
  \end{align*}
  Yet, it is not normal. If we have transpositions $x = \left( x_i \right)_{i} = \left( \left( i,i+1 \right) \right)_{i\in \N}$ and $y = \left( y_i \right)_{i} = \left( \left( 1,i \right) \right)_{i\in \N}$, then $x\in \mathcal{N}$ but $y^{-1}xy\notin \mathcal{N}$, since $y_i^{-1}x_iy_i = \left( 1,i+1 \right)$ for each $i$, and we have $d\left( y_i^{-1}x_iy_i,e \right) \rightarrow 1/2$.

  In particular, this means that the homogeneous space $ \mathcal{G}/\mathcal{N} $ admits a $\mathcal{G}$-invariant metric $d\left( x,y \right) = \lim_{\mathcal{U}}d\left( x_n,y_n \right)$, it is not a group.
\end{example}
In order to develop an ultraproduct of metric groups that is actually a group, we need to use bi-invariant metrics. That is, metrics that satisfy $d\left( x,y \right) = d\left( gx,gy \right) = d\left( xg,yg \right)$. These metrics are special because they admit topologies whose conjugation-invariant open subsets form a neighborhood basis at the identity.

If $\left( G_{\alpha},d_{\alpha} \right)_{\alpha}$ are a family of metric groups with bi-invariant metrics, then the ``null set'' subgroup $\mathcal{N}$ is normal, and for any ultrafilter $ \mathcal{U} $ on $A$, the ultraproduct
\begin{align*}
  \prod_{\mathcal{U}} G_{\alpha} &= \mathcal{G}/\mathcal{N}
\end{align*}
is indeed a group. The bi-invariant metric
\begin{align*}
  d\left( x\mathcal{N},y\mathcal{N} \right) &= \lim_{\mathcal{U}} d_{\alpha}\left( x_{\alpha},y_{\alpha} \right)
\end{align*}
yields the metric ultraproduct of the family $\left( G_{\alpha},d_{\alpha} \right)$ with respect to the ultrafilter $ \mathcal{U} $.
\begin{example}[Groups with Bi-Invariant Metrics]\hfill
  \begin{enumerate}[(i)]
    \item If $\left( X,\mu \right)$ is a finite measure space, then group $\aut\left( X,\mu \right)$ admits a bi-invariant metric given by
      \begin{align*}
        d\left( \sigma,\tau \right) &= \mu\left( \set{x\in X : \sigma(x)\neq \tau(x)} \right).
      \end{align*}
    \item The (normalized) Hamming distance on the symmetric group $S_n$ is given by
      \begin{align*}
        d_{H}\left( \sigma,\tau \right) &= \frac{1}{n}\left\vert \set{i\in [n] : \sigma(i)\neq \tau(i)} \right\vert.
      \end{align*}
      This is the special case of a finite set $[n]$ equipped with the (normalized) counting measure.
    \item If $H$ is a Hilbert space, then $U\left( H \right)$ denotes the group of unitary operators on $H$. The \textit{uniform operator metric} is given by
      \begin{align*}
        d_{u}\left( u,v \right) &= \norm{u-v}.
      \end{align*}
      This is a bi-invariant metric.
    \item If $H = \C^n$ is an $n$-dimensional Hilbert space, then $U(H) = U(n)$ is the unitary group of rank $n$. The \textit{normalized Hilbert--Schmidt metric} on $U(n)$ is given by
      \begin{align*}
        d_{\operatorname{HS}}\left( u,v \right) &= \sqrt{\tr_n\left( (u-v)^{\ast}(u-v) \right)},
      \end{align*}
      where $\tr_n$ is the normalized trace $\tr_n = \frac{1}{n}\Tr$.
  \end{enumerate}
\end{example}
Every metric group $G$ embeds into its ultrapower $ G^{\mathcal{U}} $ with the diagonal embedding.
\subsection{Approximate Homomorphisms}%
\subsection{Cayley Graphs}%

\nocite{loh_geometric_group_theory,topics_groups_geometry,cellular_automata_and_groups,pestov2008hyperlinearsoficgroupsbrief,pestov2012introductionhyperlinearsoficgroups,hayes2020randommatrixapproachabsorption}
\printbibliography
\end{document}
